% Part: normal-modal-logic
% Chapter: frame-correspondence
% Section: standard-translation

\documentclass[../../../include/open-logic-section]{subfiles}

\begin{document}

\olfileid[pt]{nml}{frd}{st}

\olsection{Definibilidade de segunda ordem}

Nem toda propriedade de frames definível por fórmulas modais é definível em
primeira ordem. Contudo, se permitirmos a quantificação sobre predicados de um
lugar (isto é, a quantificação monádica de segunda ordem), definiremos todas
as propriedades de frames modalmente definíveis. O recurso consiste em
explorar uma maneira sistemática pela qual as condições em que uma
!!{formula} modal é verdadeira em um mundo se relacionam com !!{formula}s de
primeira ordem. Trata-se da chamada tradução padrão de fórmulas modais para
fórmulas de primeira ordem em uma linguagem que contém não apenas um
!!{predicate}~$Q$ de dois lugares para a relação de acessibilidade, mas também
um !!{predicate}~$P_i$ de um lugar para as !!{propositional variable}s
~$\Obj p_i$ que aparecem em~$!A$.

\begin{defn}
  A \emph{tradução padrão}~$\ST_x(!A)$ é definida indutivamente da seguinte
  maneira:
  \begin{enumerate}
  \tagitem{prvFalse}{\indcase{!A}{\lfalse}{$\ST_x(\indfrm) = \lfalse$.}}{}
  \tagitem{prvTrue}{\indcase{!A}{\lfalse}{$\ST_x(\indfrm) = \ltrue$.}}{}
  \item \indcase{!A}{\Obj p_i}{$\ST_x(\indfrm) = \Atom{P_i}{x}$.}
  \tagitem{prvNot}{\indcase{!A}{\lnot !B}{$\ST_x(\indfrm) = \lnot \ST_x(!B)$.}}{}
  \tagitem{prvAnd}{\indcase{!A}{(!B \land !C)}{$\ST_x(\indfrm) = (\ST_x(!B)
      \land \ST_x(!C))$.}}{}
  \tagitem{prvOr}{\indcase{!A}{(!B \lor !C)}{$\ST_x(\indfrm) = (\ST_x(!B)
      \lor \ST_x(!C))$.}}{}
  \tagitem{prvIf}{\indcase{!A}{(!B \lif !C)}{$\ST_x(\indfrm) = (\ST_x(!B)
      \lif \ST_x(!C))$.}}{}
  \tagitem{prvIff}{\indcase{!A}{(!B \liff !C)}{$\ST_x(\indfrm) = (\ST_x(!B)
      \liff \ST_x(!C))}$.}{}
  \tagitem{prvBox}{\indcase{!A}{\Box !B}{$\ST_x(\indfrm) =
      \lforall[y][(\Atom{Q}{x,y} \lif \ST_y(!B))]$.}}{}
  \tagitem{prvDiamond}{\indcase{!A}{\Diamond !B}{$\ST_x(\indfrm) =
      \lexists[y][(\Atom{Q}{x,y} \land \ST_y(!B))]$.}}{}
  \end{enumerate}
\end{defn}

Por exemplo, $\ST_x(\Box p \lif p)$ é
$\lforall[y][(\Atom{Q}{x,y} \lif \Atom{P}{y})] \lif \Atom{P}{x}$. Toda
!!{structure} para a linguagem de~$\ST_x(!A)$ requer um domínio, uma relação
de dois lugares atribuída a~$Q$, e subconjuntos do domínio atribuídos aos
!!{predicate}s~$P_i$ de um lugar. Em outras palavras, os componentes de tal
!!a{structure} são exatamente os de um modelo para~$!A$: o domínio é o
conjunto de mundos, a relação de dois lugares atribuída a~$Q$ é a relação de
acessibilidade, e os subconjuntos atribuídos a~$P_i$ são precisamente as
atribuições~$V(\Obj p_i)$. Não será surpreendente que a satisfação de $!A$ em
um modelo modal coincida com a de $\ST_x(!A)$ na !!{structure}
correspondente:

\begin{prop}\ollabel{prop:st}
  Seja $\mModel{M} = \tuple{W, R, V}$, seja $\Struct{M'}$ a !!{structure} de
  primeira ordem com $\Domain{M'} = W$, $\Assign{Q}{M'} = R$ e
  $\Assign{P_i}{M'} = V(\Obj p_i)$, e seja $s(x) = w$. Então
  \[
  \mSat{M}{!A}[w] \text{ se e somente se } \Sat{M'}{\ST_x(!A)}[s]
  \]
\end{prop}

\begin{proof}
  Por indução sobre~$!A$.
\end{proof}

\begin{prop}
  Suponha que $!A$ seja uma !!{formula} modal e que
  $\mModel{F} = \tuple{W, R}$ seja um frame. Seja $\Struct{F'}$ a
  !!{structure} de primeira ordem com $\Domain{F'} = W$ e
  $\Assign{Q}{F'} = R$, e seja $!A'$ a fórmula de segunda ordem
  \[
  \lforall[X_1][\dots\lforall[X_n][\lforall[x][
        \SSubst{\ST_x(!A)}{\subst{X_1}{P_1}, \dots,
          \subst{X_n}{P_n}}]]],
  \]
  onde $P_1$, \dots, $P_n$ são todos os !!{predicate}s de um lugar
  em~$\ST_x(!A)$. Então
  \[
  \mModel{F} \Entails !A \text{ se e somente se } \Sat{F'}{!A'}
  \]
\end{prop}

\begin{proof}
  $\Sat{F'}{!A'}$ se e somente se, para toda !!{structure}~$\mModel{M'}$ na
  qual $\Assign{P_i}{M'} \subseteq W$ para $i = 1$, \dots,~$n$, e para todo
  $s$ com $s(x) \in W$, $\Sat{M'}{\ST_x(!A)}[s]$. Por \olref{prop:st}, isso
  ocorre se e somente se, para todos os modelos~$\mModel{M}$ baseados
  em~$\mModel{F}$ e todo mundo~$w \in W$, $\mSat{M}{!A}[w]$; isto é,
  $\mModel{F} \Entails !A$.
\end{proof}

\begin{defn}
  Uma classe~$\mClass{F}$ de frames é \emph{definível em segunda ordem} se
  existe !!a{sentence}~$!A$ na linguagem de segunda ordem com um único
  !!{predicate}~$P$ de dois lugares e quantificadores apenas sobre variáveis
  monádicas de conjuntos, tal que
  $\mModel{F} = \tuple{W, R} \in \mClass{F}$ se e somente se $\Sat{M}{!A}$
  na !!{structure}~$\Struct{M}$ com $\Domain{M} = W$ e
  $\Assign{P}{M} = R$.
\end{defn}

\begin{cor}
  Se uma classe de frames é definível por !!a{formula}~$!A$, a classe
  correspondente de relações de acessibilidade é definível por um enunciado
  monádico de segunda ordem.
\end{cor}

\begin{proof}
  O enunciado monádico de segunda ordem~$!A'$ da demonstração anterior possui
  a propriedade exigida.
\end{proof}

Como exemplo, considere novamente a !!{formula}~$\Box p \lif p$. Ela define a
reflexividade. Naturalmente, a reflexividade é definível em primeira ordem
pelo !!{sentence}~$\lforall[x][\Atom{Q}{x,x}]$. Mas ela também é definível
pelo !!{sentence} monádico de segunda ordem
\[
\lforall[X][\lforall[x][(\lforall[y][(\Atom{Q}{x,y} \lif \Atom{X}{y})]
    \lif \Atom{X}{x})]].
\]
Isso significa, naturalmente, que os dois enunciados são equivalentes. Eis
como você pode se convencer disso diretamente: suponha primeiro que o
enunciado de segunda ordem seja verdadeiro em !!a{structure}~$\Struct{M}$.
Como $x$ e $X$ são universalmente quantificados, o restante deve valer para
qualquer $x \in W$ e qualquer conjunto~$X \subseteq W$, por exemplo, o
conjunto $\Setabs{z}{Rxz}$, onde $R = \Assign{Q}{M}$. Assim, para todo~$s$ com
$s(x) \in W$ e $s(X) = \Setabs{z}{Rxz}$, temos
$\Sat{M}{\lforall[y][(\Atom{Q}{x,y} \lif \Atom{X}{y})] \lif
    \Atom{X}{x}}$. Mas, pela maneira como escolhemos $s(X)$, isso significa
$\Sat{M}{\lforall[y][(\Atom{Q}{x,y} \lif \Atom{Q}{x,y})] \lif
    \Atom{Q}{x,x}}[s]$, que é equivalente a $\Atom{Q}{x,x}$, pois o
antecedente é válido. Como $s(x)$ é arbitrário, temos
$\Sat{M}{\lforall[x][\Atom{Q}{x,x}]}$.

Suponha agora que $\Sat{M}{\lforall[x][\Atom{Q}{x,x}]}$ e mostre que
$\Sat{M}{\lforall[X][\lforall[x][(\lforall[y][(\Atom{Q}{x,y} \lif
        \Atom{X}{y})] \lif \Atom{X}{x})]]}$. Escolha uma atribuição~$s$
qualquer e suponha $\Sat{M}{\lforall[y][(\Atom{Q}{x,y} \lif
    \Atom{X}{y})]}[s]$. Seja $s'$ a variante em~$y$ de~$s$ com
$s'(y) = s(x)$; temos $\Sat{M}{\Atom{Q}{x,y} \lif \Atom{X}{y}}[s']$, isto é,
$\Sat{M}{\Atom{Q}{x,x} \lif \Atom{X}{x}}[s]$. Como
$\Sat{M}{\lforall[x][\Atom{Q}{x,x}]}$, o antecedente é verdadeiro, e temos
$\Sat{M}{\Atom{X}{x}}[s]$, que é o que precisávamos mostrar.

Como algumas classes de frames definíveis não são definíveis em primeira
ordem, nem todo !!{sentence} monádico de segunda ordem da forma~$!A'$ é
equivalente a um !!{sentence} de primeira ordem. Não existe um método efetivo
para decidir quais são.

\end{document}
